%! Polynomial Equations, Inequalities, and Modeling — Unit 2, Lesson 2.7 — Homework
\documentclass[10pt]{article}
\usepackage{precalculus-article}
\usepackage{precalculus-boxes}
\newcommand{\TallMath}[1]{$\displaystyle #1\rule[-1.4em]{0pt}{3.2em}$}

\begin{document}

\pageheader{Unit 2, Lesson 2.7}{Homework}

\vspace{0.12in}

\begin{enumerate}[itemsep=10pt]

\item Solve each equation by factoring. Set one side to $0$ first.

\begin{enumerate}[label=(\alph*), itemsep=6pt]
  \item $x^3 - 7x^2 + 12x = 0$

        \begin{work}
          x^3 - 7x^2 + 12x &= 0 \\
          x\big(x^2 - 7x + 12\big) &= 0 \\
          x(x-3)(x-4) &= 0
        \end{work}

        Solutions: \blank{5cm}

  \item $x^3 = 9x$

        \begin{work}
          x^3 - 9x &= 0 \\
          x\big(x^2 - 9\big) &= 0 \\
          x(x-3)(x+3) &= 0
        \end{work}

        Solutions: \blank{5cm}

        Which solution disappears if you divide both sides by $x$ in the first
        step? \quad \blank{2cm}
\end{enumerate}

\item Solve $(x+2)(x-1)(x-5) < 0$ by sign analysis.

\smallskip
\renewcommand{\arraystretch}{1.5}
\begin{center}
\begin{tabular}{|l|c|c|c|c|}
\hline
\textbf{Interval} & $(-\infty,-2)$ & $(-2,1)$ & $(1,5)$ & $(5,\infty)$ \\
\hline
\textbf{Test value} & \blank{1.2cm} & \blank{1.2cm} & \blank{1.2cm} & \blank{1.2cm} \\
\hline
\textbf{Sign} & \blank{1.2cm} & \blank{1.2cm} & \blank{1.2cm} & \blank{1.2cm} \\
\hline
\end{tabular}
\end{center}
\smallskip

Solution in interval notation: \blank{5cm}

\item Let $p(x) = (x-3)^2(x+1)$. Solve $p(x) \ge 0$.

\begin{enumerate}[label=(\alph*), itemsep=6pt]
  \item Zeros and multiplicities: \blank{6cm}
  \item Sign on $(-\infty,-1)$: \blank{1.4cm} \quad
        on $(-1,3)$: \blank{1.4cm} \quad
        on $(3,\infty)$: \blank{1.4cm}
  \item Solution in interval notation: \blank{4cm}
  \item At which zero does the sign \textit{fail} to change, and what feature
        of that zero explains it? \quad \blank{6cm}
\end{enumerate}

\boxguard[20]
\item A sheet of card measures 18 cm by 24 cm. Squares of side $x$ are cut from
      the four corners and the flaps folded up to make an open box, so
      $V(x) = x(18-2x)(24-2x)$.

\begin{enumerate}[label=(\alph*), itemsep=6pt]
  \item State the domain restriction the context forces, and say in a few words
        where it comes from. \quad \blank{5cm}
  \item Write $V(x)$ in expanded form.

        \begin{work}
          V(x) &= x(18-2x)(24-2x) \\
               &= x\big(432 - 36x - 48x + 4x^2\big) \\
               &= x\big(4x^2 - 84x + 432\big) \\
               &= 4x^3 - 84x^2 + 432x
        \end{work}

        $V(x) = $ \blank{6cm}
  \item Find $V(4)$ and report it in a sentence with units.

        \vspace{0.05in}
        \writeline
\end{enumerate}

\item A data set has equally spaced inputs. Its successive differences are the
      route to the right regression.

\smallskip
\renewcommand{\arraystretch}{1.4}
\begin{center}
\begin{tabular}{|c|c|c|c|c|c|c|}
\hline
$x$ & $0$ & $1$ & $2$ & $3$ & $4$ & $5$ \\
\hline
$y$ & $2$ & $2$ & $8$ & $26$ & $62$ & $122$ \\
\hline
\end{tabular}
\end{center}
\smallskip

\begin{enumerate}[label=(\alph*), itemsep=6pt]
  \item First differences: \blank{6cm}
  \item Second differences: \blank{5cm}
  \item Third differences: \blank{4cm}
  \item Which regression --- linear, quadratic, cubic, or quartic --- fits this
        data exactly, and how do the differences tell you?
        \quad \blank{5cm}
  \item The calculator returns the coefficients. Name one thing it does
        \textit{not} decide for you. \quad \blank{5cm}
\end{enumerate}

\item The solution set of $x(x-2)^2 > 0$ is

\begin{enumerate}[label=\Alph*., itemsep=3pt]
  \item $(0, \infty)$
  \item $(0,2) \cup (2,\infty)$
  \item $[0, \infty)$
  \item $(2, \infty)$
\end{enumerate}

\end{enumerate}

\vspace{0.1in}

\boxguard
\begin{extensionbox}
Why is one test value per interval enough? Suppose $p$ is positive at some
point of an interval and negative at another point of the same interval, and
that the interval contains no zeros of $p$. Explain, in two or three sentences,
why that is impossible. (Your argument will use the fact that a polynomial
graph has no breaks in it --- calculus will give that fact a name.)

\vspace{0.05in}
\writelines{3}
\end{extensionbox}

\vspace{0.1in}

\boxguard
\begin{notesbox}{Preview: The Unit 2 Test, and Unit 3}
That closes Unit 2. The \textbf{Unit 2 Test} runs the whole arc: end behavior,
real and complex zeros, multiplicity, equivalent forms, and today's equations,
inequalities, and models. Then \textbf{Unit 3, Rational Functions}, divides one
polynomial by another. Today's sign chart shows up again almost immediately
there --- with one new rule. A rational expression can change sign at a zero of
its \textit{denominator} too, at a point that is not even in the domain.
\end{notesbox}

\end{document}
